% GENERADO por benchmarks/np2_sequence/tesis/etapa3_tablas.jl; no editar a mano.
% Procedencia:
%   carbon_rohf_results_3P_R30.0.jld2 (sha256 e8c5884b)
%   CI carbon_rohf_results_3P_R30.0.jld2 m=40 (sha256 e8c5884b, commit 6dfb84b)
%   silicon_rohf_results_3P_R30.0.jld2 (sha256 09adeb32)
%   CI silicon_rohf_results_3P_R30.0.jld2 m=20 (sha256 09adeb32, commit 2479092)
%   germanium_rohf_results_3P_R30.0.jld2 (sha256 7567c491)
%   CI germanium_rohf_results_3P_R30.0.jld2 m=20 (sha256 7567c491, commit 2479092)
%   germanium_rohf_results_3P_R30.0_ad0.750.jld2 (sha256 1d18de80)
%   CI germanium_rohf_results_3P_R30.0_ad0.750.jld2 m=20 (sha256 1d18de80, commit 2479092) <- alpha_d = 0.75
%   tin_rohf_results_3P_R30.0.jld2 (sha256 7db46b8f)
%   CI tin_rohf_results_3P_R30.0.jld2 m=20 (sha256 7db46b8f, commit 2479092)
%   tin_rohf_results_3P_R30.0_ad3.000.jld2 (sha256 b6dcaa41)
%   CI tin_rohf_results_3P_R30.0_ad3.000.jld2 m=20 (sha256 b6dcaa41, commit 2479092) <- alpha_d = 3.00
%   En eV, con el error relativo frente al NIST entre parentesis. HF+CI = -eps_np - E_corr: el ion np^1 no tiene correlacion de pareja de valencia.
\begin{tabular}{lcccc}
    \toprule
    \textbf{Elemento} & \textbf{Koopmans} & \textbf{HF+CI} & \textbf{CI+$V_{\text{pol}}$} & \textbf{NIST} \\
    \midrule
    Carbono & 11.79 ($+4.7$\%) & 12.06 ($+7.1$\%) & -- & 11.26 \\
    Silicio & 8.08 ($-0.8$\%) & 8.32 ($+2.0$\%) & -- & 8.15 \\
    Germanio & 7.82 ($-1.0$\%) & 8.01 ($+1.4$\%) & 8.28 ($+4.8$\%) & 7.90 \\
    Estaño & 7.21 ($-1.8$\%) & 7.38 ($+0.5$\%) & 7.96 ($+8.4$\%) & 7.34 \\
    \bottomrule
\end{tabular}
